\small
\linespread{0.94}\selectfont
\setlength{\parskip}{0pt}
\setlength{\abovedisplayskip}{4pt plus 1pt minus 1pt}
\setlength{\belowdisplayskip}{4pt plus 1pt minus 1pt}
\setlength{\abovedisplayshortskip}{3pt plus 1pt minus 1pt}
\setlength{\belowdisplayshortskip}{3pt plus 1pt minus 1pt}
\setlength{\textfloatsep}{8pt plus 2pt minus 2pt}
\setlength{\floatsep}{7pt plus 2pt minus 2pt}
\setlength{\intextsep}{7pt plus 2pt minus 2pt}
\setlength{\abovecaptionskip}{3pt}
\setlength{\belowcaptionskip}{2pt}

\section{问题重述}

随着人工智能推理与训练任务持续增长，数据中心的算力分配已经与能源供给、碳排放和通信服务质量形成紧密耦合。题目给出了六个区域的 GPU 容量、IT 与设施功率边界、逐时非 AI 负荷、PUE、电价、碳强度、新能源出力、购售电限制、储能参数以及区域间单向网络时延，同时提供了三类异构 AI 任务的到达时刻、GPU 需求、预计持续时间、来源区域、最大允许时延和完成期限。需要在满足物理边界与服务约束的条件下，合理确定任务的执行区域和开工时刻，并协调新能源、储能及电网购售电行为。

整个运行过程具有明确的分层时间边界。第 $0$ 至第 $2399$ 小时构成任务到达主时域，第 $2400$ 至第 $2405$ 小时不再接收新任务，仅允许此前到达的弹性任务继续执行并完成结清；第 $2406$ 小时不得安排任何任务，只用于电力运行和储能终端状态结算。所有任务均不可拆分、不可抢占，也不得在执行过程中迁移。实时推理任务必须在到达时刻立即开工，批量推理与 AI 训练任务可以在允许范围内延迟，但其完成时点不得超过任务截止时间和系统统一终点。

对于问题一，需要按区域和任务类型统计 GPU 需求的规模、分布及时间规律，建立短期需求预测模型，并按照时间顺序完成模型训练、验证与最终测试。预测结果仅用于模型精度评价，最后 $24$ 小时的基础调度必须使用第 $2376$ 至第 $2399$ 小时的实际到达任务。调度时还需根据分钟级持续时间计算任务与各小时的真实重叠，据此分别审计瞬时 GPU 占用和小时电量，并展示任务实际起止时刻、执行区域及区域 GPU 利用情况。

对于问题二，需要使用全体实际任务和逐时电力参数重新构造碳感知调度方案。在任务唯一执行、实时任务即时开工、网络时延、GPU 容量、IT 功率、设施功率及完成时限均得到满足的基础上，系统运行成本、购电碳排放、有效网络时延和新能源利用率应被统一复算。问题二不配置储能，新能源被用于直接供能、区域外送或弃置，不足负荷由电网购电补足，且购电与售电行为必须服从区域边界及互斥要求。

对于问题三，任务调度结果不再作为决策变量，区域 AI IT 负荷被固定为附件中的基准负荷，区域设施负荷由基准 AI IT 负荷与非 AI IT 负荷之和乘以 PUE 得到。需要在储能初始状态、荷电状态递推、容量和功率边界、充放电互斥、购售电互斥以及终端能量不透支等条件下，确定逐区域逐时的充电、放电、购电、售电、新能源分配和弃电策略，并将含储能方案与无储能比较状态在成本、碳排放、峰值净购电功率和负荷爬坡方面进行比较。

对于问题四，需要把任务迁移与开工、设施负荷、新能源分配、储能运行以及购售电决策纳入同一框架，在运行成本、碳排放、网络时延、服务质量、新能源利用率和区域峰值净购电之间进行权衡。碳约束、电价机制与新能源波动等情景应从同一联合基准独立重构，并在相同任务集、容量边界、功率映射和评价口径下接受完整审计，从而辨析不同外部条件对任务迁移方向、执行窗口与能源行为的影响。
